\begin{textsample}{POS Dim 1 – 2011 – Score 31.00 – 2011/t002439.txt}  \label{ex:f1_pos_032}
\textbf{The} oceans cover some 70 percent of our \textbf{planet}.
And I think Arthur C.
Clarke probably had it right when he said that perhaps we ought to call our \textbf{planet} Planet Ocean.
And \textbf{the} oceans are hugely productive, as you can see by \textbf{the} satellite image of photosynthesis, \textbf{the} production of new life.
In fact, \textbf{the} oceans produce half of \textbf{the} new life every day on Earth as well as about half \textbf{the} oxygen that we breathe.
In addition to that, it harbors a lot of \textbf{the} biodiversity on Earth, and much of it we don ' t know about.
But I ' ll tell you some of that today.
That also doesn ' t even get into \textbf{the} whole protein extraction that we do from \textbf{the} ocean.
That ' s about 10 percent of our global needs and 100 percent of some \textbf{island} nations.
If you were to descend into \textbf{the} 95 percent of \textbf{the} biosphere that ' s livable, it would quickly become pitch black, interrupted only by pinpoints of light from bioluminescent organisms.
And if you turn \textbf{the} lights on, you might periodically see spectacular organisms swim by, because those are \textbf{the} denizens of \textbf{the} deep, \textbf{the} things that live in \textbf{the} deep ocean.
And eventually, \textbf{the} deep \textbf{sea} floor would come into view.
This type of \textbf{habitat} covers more of \textbf{the} Earth ' s \textbf{surface} than all other \textbf{habitats} combined.
And yet, we know more about \textbf{the} \textbf{surface} of \textbf{the} Moon and about Mars than we do about this \textbf{habitat}, despite \textbf{the} fact that we have yet to extract a gram of food, a breath of oxygen or a drop of water from those bodies.
And so 10 years ago, an international program began called \textbf{the} Census of Marine Life, which set out to try and improve our understanding of life in \textbf{the} global oceans.
It involved 17 different projects around \textbf{the} world.
As you can see, these are \textbf{the} footprints of \textbf{the} different projects.
And I hope you ' ll appreciate \textbf{the} level of global coverage that it managed to achieve.
It all began when two scientists, Fred Grassle and Jesse Ausubel, met in Woods Hole, Massachusetts where both were guests at \textbf{the} famed oceanographic institute.
And Fred was lamenting \textbf{the} state of marine biodiversity and \textbf{the} fact that it was in trouble and nothing was being done about it.
Well, from that discussion grew this program that involved 2, 700 scientists from more than 80 countries around \textbf{the} world who engaged in 540 ocean expeditions at a combined cost of 650 million dollars to study \textbf{the} distribution, diversity and abundance of life in \textbf{the} global ocean.
And so what did we find?
We found spectacular new species, \textbf{the} most beautiful and visually \textbf{stunning} things everywhere we looked — from \textbf{the} shoreline to \textbf{the} \textbf{abyss}, form microbes all \textbf{the} way up to \textbf{fish} and everything in between.
And \textbf{the} limiting step here wasn ' t \textbf{the} unknown diversity of life, but rather \textbf{the} taxonomic specialists who can identify and catalog these species that became \textbf{the} limiting step.
They, in fact, are an endangered species themselves.
There are actually four to five new species described everyday for \textbf{the} oceans.
And as I say, it could be a much larger number.
Now, I come from Newfoundland in Canada — It ' s an \textbf{island} off \textbf{the} east coast of that \textbf{continent} — where we experienced one of \textbf{the} worst \textbf{fishing} disasters in human history.
And so this photograph shows a small boy next to a codfish.
It ' s around 1900.
Now, when I was a boy of about his age, I would go out \textbf{fishing} with my grandfather and we would catch \textbf{fish} about half that size.
And I thought that was \textbf{the} norm, because I had never seen \textbf{fish} like this.
If you were to go out there today, 20 years after this fishery collapsed, if you could catch a \textbf{fish}, which would be a bit of a challenge, it would be half that size still.
So what we ' re experiencing is something called shifting baselines.
Our expectations of what \textbf{the} oceans can produce is something that we don ' t really appreciate because we haven ' t seen it in our lifetimes.
Now most of us, and I would say me included, think that human exploitation of \textbf{the} oceans really only became very serious in \textbf{the} last 50 to, perhaps, 100 years or so. \textbf{The} census actually tried to look back in time, using every source of information they could get their hands on.
And so anything from restaurant menus to monastery records to ships ' logs to see what \textbf{the} oceans looked like.
Because science data really goes back to, at best, World War II, for \textbf{the} most part.
And so what they found, in fact, is that exploitation really began heavily with \textbf{the} Romans.
And so at that time, of course, there was no refrigeration.
So fishermen could only catch what they could either eat or sell that day.
But \textbf{the} Romans developed salting.
And with salting, it became possible to store \textbf{fish} and to transport it long distances.
And so began industrial \textbf{fishing}.
And so these are \textbf{the} sorts of extrapolations that we have of what sort of loss we ' ve had \textbf{relative} to pre-human impacts on \textbf{the} ocean.
They range from 65 to 98 percent for these major groups of organisms, as shown in \textbf{the} dark \textbf{blue} bars.
Now for those species \textbf{the} we managed to leave alone, that we protect — for example, marine mammals in recent years and \textbf{sea} birds — there is some recovery.
So it ' s not all hopeless.
But for \textbf{the} most part, we ' ve gone from salting to exhausting.
Now this other line of evidence is a really interesting one.
It ' s from trophy \textbf{fish} caught off \textbf{the} coast of Florida.
And so this is a photograph from \textbf{the} 1950s.
I want you to notice \textbf{the} scale on \textbf{the} \textbf{slide}, because when you see \textbf{the} same picture from \textbf{the} 1980s, we see \textbf{the} \textbf{fish} are much smaller and we ' re also seeing a change in terms of \textbf{the} composition of those \textbf{fish}.
By 2007, \textbf{the} catch was actually laughable in terms of \textbf{the} size for a trophy \textbf{fish}.
But this is no laughing matter. \textbf{The} oceans have lost a lot of their productivity and we ' re responsible for it.
So what ' s left?
Actually quite a lot.
There ' s a lot of exciting things, and I ' m going to tell you a little bit about them.
And I want to start with a bit on technology, because, of course, this is a TED Conference and you want to hear something on technology.
So one of \textbf{the} tools that we use to sample \textbf{the} deep ocean are remotely operated vehicles.
So these are tethered vehicles we lower down to \textbf{the} \textbf{sea} floor where they ' re our eyes and our hands for working on \textbf{the} \textbf{sea} bottom.
So a couple of years ago, I was supposed to go on an oceanographic cruise and I couldn ' t go because of a scheduling conflict.
But through a satellite link I was able to sit at my study at home with my dog curled up at my feet, a cup of tea in my hand, and I could tell \textbf{the} pilot,"I want a sample right there."And that ' s exactly what \textbf{the} pilot did for me.
That ' s \textbf{the} sort of technology that ' s available today that really wasn ' t available even a decade ago.
So it allows us to sample these amazing \textbf{habitats} that are very far from \textbf{the} \textbf{surface} and very far from light.
And so one of \textbf{the} tools that we can use to sample \textbf{the} oceans is acoustics, or sound waves.
And \textbf{the} advantage of sound waves is that they actually pass well through water, unlike light.
And so we can send out sound waves, they bounce off \textbf{objects} like \textbf{fish} and are reflected back.
And so in this example, a census scientist took out two ships.
One would send out sound waves that would bounce back.
They would be received by a second ship, and that would give us very precise estimates, in this case, of 250 billion herring in a period of about a minute.
And that ' s an area about \textbf{the} size of Manhattan Island.
And to be able to do that is a tremendous fisheries tool, because knowing how many \textbf{fish} are there is really critical.
We can also use satellite \textbf{tags} to track animals as they move through \textbf{the} oceans.
And so for animals that come to \textbf{the} \textbf{surface} to breathe, such as this \textbf{elephant} \textbf{seal}, it ' s an opportunity to send data back to shore and tell us where exactly it is in \textbf{the} ocean.
And so from that we can produce these tracks.
For example, \textbf{the} dark \textbf{blue} shows you where \textbf{the} \textbf{elephant} \textbf{seal} moved in \textbf{the} north Pacific.
Now I realize for those of you who are colorblind, this \textbf{slide} is not very helpful, but stick with me nonetheless.
For animals that don ' t \textbf{surface}, we have something called pop-up \textbf{tags}, which collect data about light and what time \textbf{the} sun rises and sets.
And then at some period of time it pops up to \textbf{the} \textbf{surface} and, again, relays that data back to shore.
Because GPS doesn ' t work under water.
That ' s why we need these tools.
And so from this we ' re able to identify these \textbf{blue} \textbf{highways}, these hot spots in \textbf{the} ocean, that should be real priority areas for ocean conservation.
Now one of \textbf{the} other things that you may think about is that, when you go to \textbf{the} \textbf{supermarket} and you buy things, they ' re scanned.
And so there ' s a barcode on that product that tells \textbf{the} \textbf{computer} exactly what \textbf{the} product is.
Geneticists have developed a similar tool called genetic barcoding.
And what barcoding does is use a specific gene called CO1 that ' s consistent within a species, but varies among species.
And so what that means is we can unambiguously identify which species are which even if they look similar to each other, but may be biologically quite different.
Now one of \textbf{the} nicest examples I like to cite on this is \textbf{the} story of two young women, high school students in New York City, who worked with \textbf{the} census.
They went out and collected \textbf{fish} from markets and from restaurants in New York City and they barcoded it.
Well what they found was mislabeled \textbf{fish}.
So for example, they found something which was sold as tuna, which is very valuable, was in fact tilapia, which is a much less valuable \textbf{fish}.
They also found an endangered species sold as a common one.
So barcoding allows us to know what we ' re working with and also what we ' re eating. \textbf{The} Ocean Biogeographic Information System is \textbf{the} database for all \textbf{the} census data.
It ' s open access ; you can all go in and download data as you wish.
And it contains all \textbf{the} data from \textbf{the} census plus other data sets that people were willing to contribute.
And so what you can do with that is to plot \textbf{the} distribution of species and where they \textbf{occur} in \textbf{the} oceans.
What I ' ve plotted up here is \textbf{the} data that we have on hand.
This is where our sampling effort has concentrated.
Now what you can see is we ' ve sampled \textbf{the} area in \textbf{the} North Atlantic, in \textbf{the} North Sea in particular, and also \textbf{the} east coast of North America fairly well.
That ' s \textbf{the} warm colors which show a well-sampled region. \textbf{The} cold colors, \textbf{the} \textbf{blue} and \textbf{the} black, show areas where we have almost no data.
So even after a 10-year census, there are large areas that still remain unexplored.
Now there are a group of scientists living in Texas, working in \textbf{the} Gulf of Mexico who decided really as a labor of love to pull together all \textbf{the} knowledge they could about biodiversity in \textbf{the} Gulf of Mexico.
And so they put this together, a list of all \textbf{the} species, where they ' re known to \textbf{occur}, and it really seemed like a very esoteric, scientific type of exercise.
But then, of course, there was \textbf{the} Deep Horizon \textbf{oil} \textbf{spill}.
So all of a sudden, this labor of love for no obvious economic reason has become a critical piece of information in terms of how that system is going to recover, how long it will take and how \textbf{the} lawsuits and \textbf{the} multi-billion-dollar discussions that are going to happen in \textbf{the} coming years are likely to be resolved.
So what did we find?
Well, I could stand here for hours, but, of course, I ' m not allowed to do that.
But I will tell you some of my favorite discoveries from \textbf{the} census.
So one of \textbf{the} things we discovered is where are \textbf{the} hot spots of diversity?
Where do we find \textbf{the} most species of ocean life?
And what we find if we plot up \textbf{the} well-known species is this sort of a distribution.
And what we see is that for coastal \textbf{tags}, for those organisms that live near \textbf{the} shoreline, they ' re most diverse in \textbf{the} tropics.
This is something we ' ve actually known for a while, so it ' s not a real breakthrough.
What is really exciting though is that \textbf{the} oceanic \textbf{tags}, or \textbf{the} ones that live far from \textbf{the} coast, are actually more diverse at intermediate latitudes.
This is \textbf{the} sort of data, again, that managers could use if they want to prioritize areas of \textbf{the} ocean that we need to conserve.
You can do this on a global scale, but you can also do it on a regional scale.
And that ' s why biodiversity data can be so valuable.
Now while a lot of \textbf{the} species we discovered in \textbf{the} census are things that are small and hard to see, that certainly wasn ' t always \textbf{the} case.
For example, while it ' s hard to believe that a three kilogram lobster could elude scientists, it did until a few years ago when South African fishermen requested an export permit and scientists realized that this was something new to science.
Similarly this Golden V kelp collected in Alaska just below \textbf{the} low water mark is probably a new species.
Even though it ' s three meters long, it actually, again, eluded science.
Now this guy, this bigfin squid, is seven meters in length.
But to be fair, it lives in \textbf{the} deep waters of \textbf{the} Mid- Atlantic Ridge, so it was a lot harder to find.
But there ' s still potential for discovery of big and exciting things.
This particular \textbf{shrimp}, we ' ve dubbed it \textbf{the} Jurassic \textbf{shrimp}, it ' s thought to have gone extinct 50 years ago — at least it was, until \textbf{the} census discovered it was living and doing just fine off \textbf{the} coast of Australia.
And it shows that \textbf{the} ocean, because of its vastness, can hide secrets for a very long time.
So, Steven Spielberg, eat your heart out.
If we look at distributions, in fact distributions change dramatically.
And so one of \textbf{the} records that we had was this sooty shearwater, which undergoes these spectacular migrations all \textbf{the} way from New Zealand all \textbf{the} way up to Alaska and back again in \textbf{search} of endless summer as they complete their life cycles.
We also talked about \textbf{the} White Shark Cafe.
This is a location in \textbf{the} Pacific where white shark converge.
We don ' t know why they converge there, we simply don ' t know.
That ' s a question for \textbf{the} future.
One of \textbf{the} things that we ' re taught in high school is that all animals require oxygen in order to survive.
Now this little critter, it ' s only about half a millimeter in size, not terribly charismatic.
But it was only discovered in \textbf{the} early 1980s.
But \textbf{the} really interesting thing about it is that, a few years ago, census scientists discovered that this guy can thrive in oxygen-poor sediments in \textbf{the} deep Mediterranean Sea.
So now they know that, in fact, animals can live without oxygen, at least some of them, and that they can adapt to even \textbf{the} harshest of conditions.
If you were to suck all \textbf{the} water out of \textbf{the} ocean, this is what you ' d be left behind with, and that ' s \textbf{the} biomass of life on \textbf{the} \textbf{sea} floor.
Now what we see is huge biomass towards \textbf{the} poles and not much biomass in between.
We found life in \textbf{the} extremes.
And so there were new species that were found that live inside ice and help to support an ice-based food \textbf{web}.
And we also found this spectacular yeti \textbf{crab} that lives near boiling hot hydrothermal vents at Easter Island.
And this particular species really captured \textbf{the} public ' s attention.
We also found \textbf{the} deepest vents known yet — 5, 000 meters — \textbf{the} hottest vents at 407 degrees Celsius — vents in \textbf{the} South Pacific and also in \textbf{the} Arctic where none had been found before.
So even new environments are still within \textbf{the} domain of \textbf{the} discoverable.
Now in terms of \textbf{the} unknowns, there are many.
And I ' m just going to summarize just a few of them very quickly for you.
First of all, we might ask, how many \textbf{fishes} in \textbf{the} \textbf{sea}?
We actually know \textbf{the} \textbf{fishes} better than we do any other group in \textbf{the} ocean other than marine mammals.
And so we can actually extrapolate based on rates of discovery how many more species we ' re likely to discover.
And from that, we actually \textbf{calculate} that we know about 16, 500 marine species and there are probably another 1, 000 to 4, 000 left to go.
So we ' ve done pretty well.
We ' ve got about 75 percent of \textbf{the} \textbf{fish}, maybe as much as 90 percent.
But \textbf{the} \textbf{fishes}, as I say, are \textbf{the} best known.
So our level of knowledge is much less for other groups of organisms.
Now this figure is actually based on a brand new paper that ' s going to come out in \textbf{the} journal PLoS \textbf{Biology}.
And what is does is predict how many more species there are on land and in \textbf{the} ocean.
And what they found is that they think that we know of about nine percent of \textbf{the} species in \textbf{the} ocean.
That means 91 percent, even after \textbf{the} census, still remain to be discovered.
And so that turns out to be about two million species once all is said and done.
So we still have quite a lot of work to do in terms of unknowns.
Now this bacterium is part of mats that are found off \textbf{the} coast of Chile.
And these mats actually cover an area \textbf{the} size of Greece.
And so this particular bacterium is actually visible to \textbf{the} naked eye.
But you can imagine \textbf{the} biomass that represents.
But \textbf{the} really intriguing thing about \textbf{the} microbes is just how diverse they are.
A single drop of seawater could contain 160 different types of microbes.
And \textbf{the} oceans themselves are thought potentially to contain as many as a billion different types.
So that ' s really exciting.
What are they all doing out there?
We actually don ' t know. \textbf{The} most exciting thing, I would say, about this census is \textbf{the} role of global science.
And so as we see in this image of light during \textbf{the} night, there are lots of areas of \textbf{the} Earth where human development is much greater and other areas where it ' s much less, but between them we see large dark areas of relatively unexplored ocean. \textbf{The} other point I ' d like to make about this is that this ocean ' s interconnected.
Marine organisms do not care about international boundaries ; they move where they will.
And so \textbf{the} importance then of global collaboration becomes all \textbf{the} more important.
We ' ve lost a lot of paradise.
For example, these tuna that were once so abundant in \textbf{the} North Sea are now effectively gone.
There were trawls taken in \textbf{the} deep \textbf{sea} in \textbf{the} Mediterranean, which collected more garbage than they did animals.
And that ' s \textbf{the} deep \textbf{sea}, that ' s \textbf{the} environment that we consider to be among \textbf{the} most pristine \textbf{left} on Earth.
And there are a lot of other pressures.
Ocean acidification is a really big issue that people are concerned with, as well as ocean warming, and \textbf{the} effects they ' re going to have on coral reefs.
On \textbf{the} scale of decades, in our lifetimes, we ' re going to see a lot of damage to coral reefs.
And I could spend \textbf{the} rest of my time, which is getting very limited, going through this litany of concerns about \textbf{the} ocean, but I want to end on a more positive note.
And so \textbf{the} grand challenge then is to try and make sure that we preserve what ' s left, because there is still spectacular beauty.
And \textbf{the} oceans are so productive, there ' s so much going on in there that ' s of relevance to humans that we really need to, even from a selfish perspective, try to do better than we have in \textbf{the} past.
So we need to recognize those hot spots and do our best to protect them.
When we look at pictures like this, they take our breath away, in addition to helping to give us breath by \textbf{the} oxygen that \textbf{the} oceans provide.
Census scientists worked in \textbf{the} rain, they worked in \textbf{the} cold, they worked under water and they worked above water trying to illuminate \textbf{the} wondrous discovery, \textbf{the} still vast unknown, \textbf{the} spectacular adaptations that we see in ocean life.
So whether you ' re a yak herder living in \textbf{the} mountains of Chile, whether you ' re a stockbroker in New York City or whether you ' re a TEDster living in Edinburgh, \textbf{the} oceans matter.
And as \textbf{the} oceans go so shall we.
Thanks for listening.

% matched lemmas: abyss, biology, blue, calculate, computer, continent, crab, elephant, fish, fishing, habitat, highway, island, left, object, occur, oil, planet, relative, sea, seal, search, shrimp, slide, spill, stunning, supermarket, surface, tag, the, web
\end{textsample}
